\subsection{A bounded restricted optimizer}
The historical practical method fits Jaccard affinities, obtains an
average-linkage order and groups consecutive leaves into compact blocks of
at most eight. Full-wave validation retains either the incumbent or the
learned compact layout. The new optional method refines \emph{maximal existing
subtrees} containing at most $s$ leaves. Each selected subtree is disjoint
from the others. It keeps the order and that subtree's leaf interval, solves
Equation~\eqref{eq:dp} on projected training waves, and accepts a replacement
only if the complete integer training objective strictly decreases.
All larger ancestors remain unchanged; ties retain the incumbent subtree.

\begin{proposition}[Restricted exactness and bounded refinement]\label{prop:jrefine}
For a fixed order and fixed disjoint refinement intervals, independent exact
fitting minimizes training cost over trees obtained by replacing only those
subtrees. Retaining the incumbent on ties gives nonincreasing integer
training cost. For intervals of sizes $m_i\le s$, the DP work is at most
$O(bns^2)$ with $O(bs^2)$ peak DP storage, apart from input projection,
cost computation and ordering.
\end{proposition}
\begin{proof}
All ancestor contributions depend only on the unchanged coverage and arity.
The mutable objective is the sum of independent interval objectives, so
minimizing each minimizes the restricted sum. The incumbent is feasible in
every subproblem. Since $\sum_im_i\le n$,
$\sum_im_i^3\le s^2\sum_im_i\le ns^2$. Sequential solution gives the stated
storage bound. Integer rescoring rejects a numerically nonimproving DP
replacement.
\end{proof}

The implementation uses float64 inside the existing recurrence; exactness
is to numerical precision. Independent integer rescoring establishes
monotonicity, not exact rational optimality for every large instance.
Independent literal-tree enumeration checks tiny cases. Every leaf remains
present exactly once and all packet arities remain feasible.

The primary frozen cap is $s=32$; $64$ is a secondary sensitivity.
A full compact8 ancestor above eight eight-leaf blocks covers $8^2=64$
leaves. With $s=32$, its blocks are solved separately: no replacement can
cross a block boundary, merge adjacent blocks or change that ancestor.
Cap64 opens this first cross-block topology decision. Ragged boundary
subtrees can have other sizes; the actual largest solved interval is reported.
This interaction explains why cap32 often ties while the prespecified cap64
sensitivity improves. It does not change the frozen primary endpoint or
isolate a higher-order causal effect.
Dense pairwise ordering still has quadratic storage and approximately
$O(n^2\log n)$ heap work. It can dominate bounded refinement and is admitted
separately under the resource budget.
